\documentclass[10pt,a4paper]{amsart}
\usepackage[margin=19mm]{geometry}
\usepackage{amsmath,amssymb,mathtools}
\usepackage[scr=rsfs,cal=euler]{mathalfa}
\usepackage{xfrac}
\usepackage[unicode,hidelinks,bookmarks=false]{hyperref}
\newcommand{\ie}{i.e.\ }
\newcommand{\cf}{cf.\ }
\newcommand{\resp}{respectively\ }
\newcommand{\Subsection}{\subsection}
\providecommand{\nobreakdash}{\nobreak\hbox{-}}
\setlength{\emergencystretch}{2em}
\multlinegap0pt
\usepackage{amssymb}
\usepackage{enumerate}
\usepackage{mathtools}
\usepackage{float}
\usepackage{xcolor}
\newtheorem{thm}{Theorem}
\newtheorem{lem}{Lemma}
\newcommand{\s}{\mathbb{S}}
\title{On Round Surgery on framed links and their Diagrams for 3-manifolds}
\author{Prerak Deep, Dheeraj Kulkarni}
\date{Source: arXiv:2501.09518v3 (CC BY 4.0)}
\begin{document}
\maketitle

\noindent\emph{Source and licence.} On Round Surgery on framed links and their Diagrams for 3-manifolds. Version arXiv:2501.09518v3 (CC BY 4.0), distributed by arXiv under the Creative Commons Attribution 4.0 International licence, http://creativecommons.org/licenses/by/4.0/, as recorded in the arXiv licence metadata for this exact version and shown on the abstract page https://arxiv.org/abs/2501.09518v3. Full licence notice: \texttt{licenses/arxiv-2501.09518-CC-BY-4.0.txt}.\par\smallskip

\noindent\textit{Editorial scope.} Counted unit: the article's lem lem (source0.tex lines 1110--1112). The article's complete original proof of it is delivered with it (source0.tex lines 1113--1211, one \texttt{proof} environment of 99 lines). No mathematics is written, repaired or re-derived: the statement and the proof are the article's own lines, in the article's own notation, and no retained line is substituted or re-flowed.  Retained uncounted as the source-local context the proof needs: lines 111--117; lines 710--712; lines 713--718; lines 775--780.  One declared adaptation: exactly 10 ancillary strings inside the retained spans are omitted (image-inclusion commands and, where the source repeats a label, the repeated label definitions) and no other line differs from the source; the figure environments, with their placement, captions and labels, are kept, and the package records the omitted strings in fidelity receipts bound to the affected bodies.  No retained line contains a citation, so the wrapper carries no bibliography and no bibliography entry.  The retained lines contain the article's own diagram code, which the wrapper typesets with the standard tikz packages; no image file is delivered and the package declares no asset dependency.  Editorial formatting only, in addition: an AMS \texttt{amsart} preamble that restates the article's own declarations and macros used by the retained lines, namely the environments \texttt{thm}, \texttt{lem} on separate counters, together with the standard packages those lines need.  The retained environments print in this document as Theorem 1, Lemma 1 in order of appearance, whereas the article prints the same items with the numbering of its own full document.  The complete native source of the article as distributed in the arXiv e-print for this exact version is bundled unchanged as \texttt{sources/171681/source0.tex}.
\begin{thm}[The Bridge Theorem]\label{BridgeTheorem}
The following two statements establish a bridge between integral Dehn surgery diagrams and round surgery diagrams:
\begin{enumerate}
    \item[(a)] Let $L=\bigcup_{i=1}^{n} (L_{i1} \cup L_{i2})$ be a round surgery diagram of joint pairs in $\mathbb{S}^3$ for a closed, connected, oriented smooth 3-manifold $M_{L}$, where $n_{ij}\in \mathbb{Z}$ is the round 1-surgery coefficient on $L_{ij}$ and $m_i\in \mathbb{Z}$ is the round 2-surgery coefficient on $L_{i2}$ for $1\leq i\leq n$ and $1 \leq j \leq 2$. Then there exists a corresponding Dehn surgery diagram on $L$ where $n_{i1}-n_{i2}+m_i$ and $m_i$ are the surgery coefficients on $L_{i1}$ and $L_{i2}$, respectively.
    \item[(b)] Conversely, given an integral Dehn surgery diagram $L= L_1 \cup \cdots \cup L_n$ of a closed, connected, oriented smooth 3-manifold with surgery coefficients $n_i \in \mathbb{Z}$ on each component $L_i$, there exists an equivalent round surgery diagram of joint pairs.
\end{enumerate}
\end{thm}

\begin{lem}[{\bf Joint pair to a Dehn surgery diagram}]\label{lem: JPtoOSD}
    Let $L=L_{11} \cup L_{12}$ be a round surgery diagram of a closed connected oriented 3-manifold such that $L$ is a joint pair of round surgeries of indices 1 and 2. Then, it corresponds to the Dehn surgery link $L$ with the surgery coefficient shown in Figure \ref{fig:JPtoOSD}.
\end{lem}

\begin{figure}[ht]
        \centering
        
        \caption{A joint pair to a Dehn surgery diagram.}
        \label{fig:JPtoOSD}
    \end{figure}

\begin{lem}[{Ordinary surgery diagram to a joint pair of round surgeries}]\label{lem: OSDtoJPRS}
Suppose $L=L_1 \cup \cdots \cup L_n$ is a Dehn surgery diagram such that component $L_i$ has a surgery coefficient $n_i$ for $1 \leq i\leq n$. 
If $n$ is even, say $n =2m$, then $(L_1 \cup L_2)\cup  \cdots \cup (L_{2m-1} \cup L_{2m})$ is a round surgery diagram of joint pairs such that each pair $L_{2i-1} \cup L_{2i}$ in the parentheses is a joint pair with round 1-surgery coefficients $n_i -n_{i+1}+k $ on $L_{i}$ and $k$ on $L_{i+1}$, and round 2-surgery coefficient $n_{i+1}$ on $L_{i+1}$, for $1 \leq i \leq m$ and $k\in \mathbb{Z}$. 

If $n$ is odd, say $n=2m-1$, then we add an unlinked unknot $L_{2m}$ with framing $(\pm1)$ to convert $L$ into a new Dehn surgery diagram with an even number of components and apply the previous statement.
\end{lem}

\begin{lem}
    The application of equivalence moves 1-4 on a given round surgery diagram $L$ does not change the resulting 3-manifold. 
\end{lem}

\begin{proof}
Suppose $L= \bigcup_{i=1}^{n} (L_{i1}\cup L_{i2})$ is a round surgery diagram of joint pairs with round 1-surgery coefficient $n_{ij}$ on component $L_{ij}$ and round 2-surgery coefficient $m_i$ on $L_{i2}$ for $1\leq i\leq n$ and $1\leq j\leq 2$. 
Below, we see that the effect of each equivalence move on $L$ does not change the resultant 3-manifold. 

{\bf Equivalence Move 1.} We apply equivalence move 1 on a joint pair of the given round surgery diagram $L$. 
Without loss of generality, we apply the equivalence move 1 on the first joint pair $L_{11} \cup L_{12}$ and obtain new round surgery coefficients on $L_{11}\cup L_{12}$.
In particular, we get a new joint pair. 
Both of the joint pairs can be converted to the same Dehn surgery diagram $D$ by Lemma \ref{lem: JPtoOSD} as shown in Figure \ref{fig: EM1Proof}. 
\begin{figure}[ht]
        \centering
        
        \caption{Diagrams involved in the equivalence move 1, i.e., diagrams on the left and right, correspond to the same Dehn surgery diagram in the middle.}
        \label{fig: EM1Proof}
\end{figure}


{\bf Equivalence Move 2.} The equivalence move 2 has two types of separate moves: The shuffle move of type A and type B. 
Since we apply a shuffle move of a certain type once at a time, we discuss their effects separately. 

{\it Shuffle Move of type A.}  Since the shuffle move of type A applies to a joint pair, without loss of generality, we apply the shuffle move A on the first joint pair $L_{11} \cup L_{12}$ of the given round surgery diagram $L$. 
\begin{figure}[ht]
        \centering
        
        \caption{Diagrams involved in the shuffle move of type A, i.e., diagrams on the left and right, correspond to the same Dehn surgery diagram in the middle.}
        \label{fig:SMTAProof}
\end{figure}
Suppose $L^{\prime}_{11} \cup L_{12}^{\prime}$ is the new joint pair obtained after performing a shuffle move of type A on $L_{11} \cup L_{12}$. 
By Theorem \ref{BridgeTheorem}, the Dehn surgery diagram corresponding to both joint pairs is the same (see Figure~\ref{fig:SMTAProof}).
Thus, the resultant 3-manifold remains unchanged. 


{\it Shuffle move of type B.} The shuffle move of type B applies to two joint pairs. In the given round surgery diagram $L$, we assume that we apply the type B move on the first two joint pairs $(L_{11} \cup L_{12}) \cup (L_{21} \cup L_{22})=: L_1 \cup L_2$.
After performing a shuffle move of type B on $L_1 \cup L_2$, we get a new round surgery diagram $L^{\prime}= \bigcup_{i=1}^{2} (L^{\prime}_{i1}\cup L^{\prime}_{12}) \bigcup_{i=3}^{n} (L_{i1}\cup L_{i2})$. We denote the new joint pairs by $L^{\prime}_i:= (L^{\prime}_{i1}\cup L^{\prime}_{12})$ for each $i=1,2$. 
\begin{figure}[ht]
    
    \caption{The shuffle move of type B applies to the diagram on the left to obtain the diagram on the right, and both of the diagrams correspond to the same surgery diagram in the middle.}
    \label{fig:SMTBProof}
\end{figure}
By Theorem \ref{BridgeTheorem}, we may convert both round surgery diagrams $L$ and $L^{\prime}$ to the same Dehn surgery diagram as shown in Figure \ref{fig:SMTBProof}. 
Thus, the shuffle move of type B does not change the resultant 3-manifold. 


{\bf Equivalence Move 3.}
This move corresponds to the addition or deletion of a joint pair $U= U_1 \cup U_2$ with round 1-surgery coefficient $k, k\pm2$ on  $U_1$ and $k$ on $U_2$, for some $k\in \mathbb{Z}$, and round 2-surgery coefficient $\pm1$ on $U_2$ to the given round surgery diagram $L$.
In particular, we get a new round surgery diagram $L^{\prime}= L \bigsqcup U$. 
By Lemma \ref{lem: OSDtoJPRS}, the joint pair $U$ corresponds to a pair of unknots $U_1 \cup U_2$ with each having Dehn surgery coefficients $n \in \{\pm1\}$. 
Recall that any unlinked disjoint $\pm1$ framed unknot corresponds to a connected sum of the resultant manifold with $\s^3$. In particular, the resultant 3-manifold is unchanged after an addition or deletion of $U$. 

{\bf Equivalence Move 4.}
We have given six types of equivalence moves for a round surgery diagram of joint pairs. 
\begin{figure}[ht]
    \centering
    
    \caption{On the left column, the round surgery diagrams are in equivalence move 4, where $L_{11}$ is replaced with $L_{11}\#_b L_{12}(m_1)$, and on the right, their corresponding surgery diagrams in Kirby move 2.}
    \label{fig:EM4-1112-Proof}
\end{figure}

\begin{figure}[ht]
    \centering
    
    \caption{On the left column, the round surgery diagrams are in equivalence move 4, where $L_{12}$ replaces with $L_{12}\#_b L_{11}(f)$, and on the right, their corresponding surgery diagrams in Kirby move 2, where framing $f= n_{11}-n_{12}+m_1$.}
    \label{fig:EM4-1211-Proof}
\end{figure}
Each of the moves does not change the resultant 3-manifold. 
For each type of move, we convert the left and right sides of the equivalence move 4 diagrams to the Dehn surgery diagrams. 
We realise that the effect of the equivalence move 4 on the round surgery diagram is equivalent to the effect of Kirby move 2 on the corresponding Dehn surgery diagrams. 
\begin{figure}[ht]
    \centering
    
    \caption{On the left column, the round surgery diagrams are in equivalence move 4, where $L_{11}$ replaces with $L_{11}\#_b L_{21}(f)$, and on the right, their corresponding surgery diagrams in Kirby move 2, where $f= n_{21}-n_{22}+m_2$}
    \label{fig:EM4-1121-Proof}
\end{figure}

\begin{figure}[ht]
    \centering
    
    \caption{On the left column, the round surgery diagrams are in equivalence move 4, where $L_{11}$ is replaced with $L_{11}\#_b L_{22}(m_2)$, and on the right, their corresponding surgery diagrams in Kirby move 2.}
    \label{fig:EM4-1122-Proof}
\end{figure}

In particular, in Figures \ref{fig:EM4-1112-Proof} and \ref{fig:EM4-1211-Proof}, we realise equivalence move 4 on a single joint pair as performing Kirby move 2 on the corresponding Dehn's surgery diagrams.
And in Figures \ref{fig:EM4-1121-Proof}, \ref{fig:EM4-1122-Proof}, \ref{fig:EM4-1221-Proof} and \ref{fig:EM4-1222-Proof}, we realise the equivalence move 4 defined on two joint pairs as performing Kirby move 2 on the corresponding Dehn's surgery diagrams. 

We know the Kirby move 2 does not change the resultant 3-manifold. Hence, the equivalence move 4 also does not change the resultant 3-manifold. 

\begin{figure}[ht]
    \centering
    
    \caption{On the left column, the round surgery diagrams are in equivalence move 4, where $L_{12}$ replaces with $L_{12}\#_b L_{21}(f)$, and on the right, their corresponding surgery diagrams in Kirby move 2, where $f=n_{21}-n_{22}+m_2$.}
    \label{fig:EM4-1221-Proof}
\end{figure}

\begin{figure}[ht]
    \centering
    
    \caption{On the left column, the round surgery diagrams are in equivalence move 4, where $L_{12}$ is replaced with $L_{12}\#_b L_{22}(m_2)$, and on the right, their corresponding surgery diagrams in Kirby move 2.}
    \label{fig:EM4-1222-Proof}
\end{figure}
\end{proof}

\end{document}
